\documentclass[12pt,a4paper]{article}

% ---- Encoding & Font ----
\usepackage{iftex}
\ifPDFTeX
  \usepackage[utf8]{inputenc}
  \usepackage[T1]{fontenc}
\else
  \usepackage{fontspec}  % XeTeX/LuaTeX (tectonic): Unicode nativo (·, —, ¿, ª, §)
\fi
\usepackage[spanish]{babel}
\usepackage{csquotes}

% ---- Page layout ----
\usepackage[top=2.5cm, bottom=2.5cm, left=2.5cm, right=2.5cm]{geometry}
\usepackage{setspace}
\onehalfspacing

% ---- Math ----
\usepackage{amsmath, amssymb, amsthm}
\usepackage{mathtools}
\usepackage{physics}
\usepackage{esint}  % \oiint

% ---- Graphics ----
\usepackage{graphicx}
\usepackage{tikz}
\usepackage{float}
\usepackage{caption}

% ---- Tables ----
\usepackage{array}
\usepackage{booktabs}
\usepackage{tabularx}
\usepackage{colortbl}

% ---- Colours ----
\usepackage{xcolor}
\definecolor{notebg}{HTML}{FFF3CD}
\definecolor{noteborder}{HTML}{FFC107}
\definecolor{boxred}{HTML}{F8D7DA}
\definecolor{boxredborder}{HTML}{F5C6CB}

% ---- Boxes ----
\usepackage{mdframed}
\usepackage{tcolorbox}
\tcbuselibrary{most}

\newtcolorbox{keybox}[1][]{
  colback=blue!5,
  colframe=blue!70!black,
  arc=4pt,
  boxrule=1pt,
  fonttitle=\bfseries,
  title={#1}
}

\newtcolorbox{warnbox}[1][]{
  colback=boxred,
  colframe=boxredborder,
  coltitle=black!75,
  arc=4pt,
  boxrule=1pt,
  fonttitle=\bfseries,
  title={#1}
}

\newtcolorbox{notebox}[1][]{
  colback=notebg,
  colframe=noteborder,
  coltitle=black!75,
  arc=4pt,
  boxrule=1pt,
  fonttitle=\bfseries,
  title={#1}
}

% ---- Hyperlinks & TOC ----
\usepackage[hidelinks]{hyperref}
\usepackage{bookmark}

% ---- Enumerate ----
\usepackage{enumitem}

% ---- Miscellaneous ----
\usepackage{icomma}
\usepackage{siunitx}
\usepackage{textcomp}
\usepackage[bottom]{footmisc}

% ---- Title ----
\title{\textbf{Clase 18: El Toroide, la Ley de Inducción de Faraday \\
        y la Ley de Lenz}\\[0.3em]
        \large{Cierre de Ampère, flujo magnético, FEM inducida y sus signos}}
\author{Transcripto y expandido --- Curso Física III (OpenFING)}
\date{}

\begin{document}

\maketitle
\thispagestyle{empty}
\newpage

\tableofcontents
\newpage

% ===================================================================
\section{Último ejemplo de Ampère: el toroide}
% ===================================================================

Cuarto ejemplo de la ley de Ampère, no calculado antes con Biot--Savart: el
\textbf{toroide}, una bobina en forma de "rosca" con radio interno $A$, radio
externo $B$ y corriente $I$. Se busca $\mathbf{B}$ entre los rollos.

\subsection{Simetría y curva amperiana}

El sistema es invariante ante rotaciones alrededor de su eje. Como las líneas
de $\mathbf{B}$ son \textbf{cerradas} (sin monopolos), son circunferencias
concéntricas: $\mathbf{B}$ tangente y de módulo constante sobre la curva
amperiana $C$ de radio $R$. Entonces
\[
\oint_C \mathbf{B}\cdot d\mathbf{l} = B\,(2\pi R).
\]

\subsection{Corriente encerrada y resultado}

Cada una de las $N$ espiras cruza la superficie una vez, en el mismo sentido:
$I_{\text{enc}} = N I$.

\begin{keybox}[Campo dentro de un toroide]
  \[
  \boxed{B = \frac{\mu_0\, N\, I}{2\pi R}}, \qquad A < R < B.
  \]
\end{keybox}

No es uniforme: decrece como $1/R$. Se obtuvo sin integrales, por simetría.

\subsection{Campo nulo fuera del toroide}

\begin{itemize}[nosep]
  \item $R > B$: toda la corriente que entra sale; neta encerrada $=0
        \Rightarrow B=0$.
  \item $R < A$: no pasa corriente $\Rightarrow B=0$.
\end{itemize}
El campo del toroide está \textbf{confinado} en su interior.

% ===================================================================
\section{La ley de inducción de Faraday}
% ===================================================================

Este es un curso de \textbf{electromagnetismo}: ya vimos que las corrientes
generan campos magnéticos; ahora veremos que \textbf{situaciones magnéticas
variables generan corrientes eléctricas}.

\subsection{Los experimentos de Faraday}

\textbf{Primer experimento (imán y espira).} Imán quieto: nada. Imán en
movimiento: aparece corriente, que cambia de signo según el sentido y
desaparece al detenerse. \emph{No es la presencia de $\mathbf{B}$, sino su
variación.} Variantes: imán mayor $\to$ efecto mayor; invertir polaridad $\to$
corriente invertida; deformar la espira (área) $\to$ efecto $\propto$ área y
$dB/dt$; girar la espira $\to$ efecto $\propto \cos\theta$ (normal contra
$\mathbf{B}$).

\textbf{Segundo experimento (electroimán).} Al cerrar un interruptor la
corriente crece y se induce; estabilizada, no hay inducción; al abrir, se
induce con signo opuesto. Da igual la fuente del campo.

\textbf{Variante clave.} Con el imán fijo y moviendo el circuito, también hay
corriente: importa la variación del \textbf{flujo}, no la de $\mathbf{B}$.

\subsection{El flujo magnético}

\begin{keybox}[Flujo magnético]
  \[
  \Phi_B = \iint_S \mathbf{B}\cdot \hat{n}\, dS.
  \]
\end{keybox}

\begin{notebox}[Superficie abierta]
  A diferencia de Gauss (superficie cerrada), aquí $S$ es abierta y su
  \textbf{borde} es la espira $C$ (cualquier superficie apoyada en el aro).
\end{notebox}

\subsection{La FEM inducida (valor absoluto)}

\begin{keybox}[Ley de Faraday --- módulo]
  \[
  \boxed{|\varepsilon| = \left|\frac{d\Phi_B}{dt}\right|}
  \]
\end{keybox}

El flujo engloba área, módulo de $\mathbf{B}$ y ángulo. La FEM existe esté o no
cerrado el circuito (si está abierto, no circula corriente). Con $N$ vueltas
iguales se multiplica por $N$.

\subsection{¿Por qué el coeficiente es 1?}

Por elección del \textbf{Sistema Internacional}: las unidades se definen para
que el coeficiente sea 1 (como el $4\pi$ de Coulomb, puesto para que no
aparezca en Gauss).

% ===================================================================
\section{Ejemplo: solenoide grande y bobina chica}
% ===================================================================

Solenoide ideal grande ($n$ espiras/longitud, corriente $I$) con una bobina
chica coaxial de $N$ vueltas y sección $A$ dentro. Se la toma "chica" para
despreciar la inducción mutua inversa. Campo dentro: $B = \mu_0 n I$. Flujo
total en la bobina chica:
\[
\Phi_B = N\,(B\,A) = \mu_0\, N\, n\, I\, A.
\]
Se impone una \textbf{rampa} de corriente:
\[
I(t) = \begin{cases} 0, & t<0 \\ \dfrac{I_{\max}}{t_f}\,t, & 0<t<t_f \\
I_{\max}, & t>t_f \end{cases}
\]
Como todo es constante salvo $I$:
\[
|\varepsilon| = \mu_0\, N\, n\, A\,\left|\frac{dI}{dt}\right|.
\]

\begin{keybox}[FEM inducida en la bobina chica]
  \[
  \boxed{|\varepsilon| = \begin{cases} 0, & t<0 \ \text{o}\ t>t_f \\[1mm]
  \dfrac{\mu_0\, N\, n\, A\, I_{\max}}{t_f}, & 0<t<t_f \end{cases}}
  \]
\end{keybox}

Reproduce a Faraday: hay FEM solo mientras la corriente cambia.

% ===================================================================
\section{Los signos: la ley de Lenz}
% ===================================================================

\subsection{Enunciado de Lenz}

Imaginando la espira cerrada (real, o por una resistencia enorme como el aire):

\begin{keybox}[Ley de Lenz]
  El signo de la corriente inducida es tal que el campo magnético que ella
  genera \textbf{se opone a la variación de flujo} que la provocó.
\end{keybox}

\begin{warnbox}[Dos matices]
  (1) Se opone a la \textbf{variación} del flujo, no al campo. (2) La oposición
  \textbf{no alcanza} a compensar del todo (no es inducción perfecta). Si
  alimentara la causa, la corriente crecería sin límite.
\end{warnbox}

\subsection{Ejemplos con imán y espira}

Con convención fija de $I$ y $\hat n$:

\begin{center}
\begin{tabular}{@{}llll@{}}
\toprule
Caso & $\Phi_B$ & ¿Crece/decrece? & $I$ inducida \\
\midrule
Imán acercándose, $\mathbf{B}\parallel\hat n$ & $+$ & crece & negativa \\
Imán alejándose, $\mathbf{B}\parallel\hat n$ & $+$ & decrece & positiva \\
Imán invertido, acercándose & $-$ & decrece & positiva \\
\bottomrule
\end{tabular}
\end{center}

\textbf{Método:} (1) flujo con su signo; (2) ¿crece o decrece?; (3) corriente
que se opone a la \emph{variación}.

\begin{notebox}[Test de consistencia]
  Al invertir el imán, la corriente debe cambiar de sentido.
\end{notebox}

\subsection{Ley de Faraday con signos}

Ligando las convenciones de FEM y flujo por la regla de la mano derecha
(orientación de $C$ $\to$ normal a $S$):

\begin{keybox}[Ley de Faraday con signo]
  \[
  \boxed{\varepsilon = -\frac{d\Phi_B}{dt}}
  \]
\end{keybox}

El signo menos codifica Lenz. Verificado en los ejemplos: flujo creciente
($d\Phi/dt>0$) $\to$ FEM negativa; flujo decreciente $\to$ FEM positiva.

% ===================================================================
\section{Ejemplo: espira que sale de un campo (frenos magnéticos)}
% ===================================================================

Campo uniforme entrante $B$; espira rectangular (lado vertical $L$, horizontal
$D$) con porción $x$ dentro, moviéndose hacia afuera con velocidad $v$;
resistencia $R$. Flujo (normal entrante): $\Phi_B = B\,D\,x$. Como
$dx/dt = -v$:
\[
\varepsilon = -\frac{d\Phi_B}{dt} = B\,D\,v.
\]

\begin{keybox}[Corriente inducida]
  \[
  \boxed{I = \frac{\varepsilon}{R} = \frac{B\,D\,v}{R}}
  \]
\end{keybox}

\subsection{Balance de energía y frenos magnéticos}

Potencia disipada: $P_{\text{dis}} = R I^2 = B^2 D^2 v^2/R$. Las fuerzas
$\mathbf{F}=I\mathbf{L}\times\mathbf{B}$ sobre los lados horizontales se
cancelan; la del lado interior se opone al movimiento:
\[
F_2 = I\,D\,B = \frac{B^2 D^2 v}{R}.
\]
La fuerza externa que mantiene $v$ constante entrega
\[
P_{\text{ext}} = F_{\text{ext}}\,v = \frac{B^2 D^2 v^2}{R} = P_{\text{dis}}.
\]

\begin{notebox}[Frenos magnéticos]
  La potencia entregada iguala a la disipada (test de consistencia). Sin fuerza
  externa, la fuerza magnética frena el sistema: es el principio de los frenos
  magnéticos (sin contacto, no se gastan; la energía se pierde por Joule).
\end{notebox}

\begin{notebox}[Próxima clase]
  Corrientes parásitas, generadores y motores, campos eléctricos inducidos.
\end{notebox}

\end{document}
